\section{Stopping and convergence}
\label{sec:rates}
The step uses the observed gap and does not require $\gamma$.
The known value of $\gamma$ is needed for the lower bound in
\eqref{eq:stop}. Quasar-convexity gives
\begin{equation}\label{eq:retained_bound}
 0\le f(x_k)-f(x^*)\le f(x_k)-\ell_k\le G(x_k)/\gamma.
\end{equation}

\begin{theorem}[Finite certified stopping]\label{thm:gap_stopping}
Under \eqref{eq:gap_smoothness} and \eqref{eq:quasar}, the algorithm
stops with $f(x_N)-f(x^*)\le\varepsilon$.
If $G(x_0)>\gamma\varepsilon$, put
$J=\lceil\log_2(G(x_0)/(\gamma\varepsilon))\rceil$.
The number of updates satisfies
\begin{equation}\label{eq:gap_iteration_bound}
 N\le\left\lceil
 \frac{\max\{4,4L_1D^2/\log 2\}}{\gamma}\,J
 +\frac{16L_0D^2}{(\log 2)\gamma^2\varepsilon}
 \right\rceil.
\end{equation}
There are exactly $N+1$ objective evaluations, $N+1$ gradient
evaluations, and $N+1$ full linear minimizations, including the terminal
query. If $G(x_0)\le\gamma\varepsilon$, the algorithm stops immediately.
\end{theorem}
\begin{proof}
The affine case $L_0=L_1=0$ takes at most one full step.
Otherwise \eqref{eq:step_lower} and \eqref{eq:descent} imply
\begin{equation}\label{eq:gap_progress}
 f(x_k)-f(x_{k+1})\ge
 \min\left\{\frac{G(x_k)}2,
 \frac{(\log 2)G(x_k)^2}{2D^2(L_0+L_1G(x_k))}\right\}.
\end{equation}
The objective decreases and $\ell_k$ increases, so $f(x_k)-\ell_k$
decreases. Set $b_j=G(x_0)/(\gamma2^j)$ for $0\le j<J$ and assign
each nonterminal update to $b_j/2<f(x_k)-\ell_k\le b_j$.
In this phase, \eqref{eq:retained_bound} gives $G(x_k)>\gamma b_j/2$.
Since $t^2/(L_0+L_1t)$ increases for $t>0$, each such update decreases
the objective by at least
\[
 \min\left\{\frac{\gamma b_j}{4},
 \frac{(\log 2)\gamma^2b_j^2}{8D^2(L_0+L_1\gamma b_j/2)}\right\}.
\]
The total decrease from updates starting in the phase is at most the
true error at its entry, hence at most $b_j$. This includes the last
update leaving the phase, whose endpoint is still feasible.
The number of updates in the phase is therefore at most
\[
 \max\left\{\frac4\gamma,
 \frac{8L_0D^2}{(\log 2)\gamma^2b_j}
 +\frac{4L_1D^2}{(\log 2)\gamma}\right\}.
\]
Summation and $\sum_{j<J}b_j^{-1}<2/\varepsilon$ prove
\eqref{eq:gap_iteration_bound}. The same finite bound holds for every
unfinished prefix, excluding an infinite nonterminal run.
At each visited point the three oracles are called once; the explicit
step \eqref{eq:log_step} makes no trial objective calls.
\end{proof}

\subsection{The regime \texorpdfstring{$L_0=0$}{L0=0}}
\begin{corollary}[Geometric convergence]\label{cor:pure_gap_linear}
If $L_0=0<L_1$, set
\[
 \underline\alpha_0=\min\{1,(\log 2)/(L_1D^2)\}.
\]
For every computed iterate,
\begin{align*}
 f(x_k)-f(x^*)&\le
 (1-\gamma\underline\alpha_0/2)^k\bigl(f(x_0)-f(x^*)\bigr),\\
 f(x_k)-\ell_k&\le
 (1-\gamma\underline\alpha_0/2)^k G(x_0)/\gamma.
\end{align*}
If the initial test fails, stopping takes at most
\[
 \left\lceil\frac{\log(G(x_0)/(\gamma\varepsilon))}
 {-\log(1-\gamma\underline\alpha_0/2)}\right\rceil
\]
updates.
\end{corollary}
\begin{proof}
The step bound gives $\alpha_k\ge\underline\alpha_0$, so decrease
is at least $\underline\alpha_0G(x_k)/2$.
Use \eqref{eq:quasar} for the first recurrence and
\eqref{eq:retained_bound}, together with $\ell_{k+1}\ge\ell_k$,
for the second.
\end{proof}
Here $\nabla^2f(x^*)=0$ follows from $G(x^*)=0$.
The nonconvex example in Section~\ref{sec:examples} has this property
and a nonzero gradient at the minimum.

\subsection{Strongly convex feasible sets}
Use the convention of \cite{vyguzov}: $Q$ is strongly convex with
constant $\lambda>0$ if the closed ball of radius
$\lambda\|u-v\|^2$ about $(u+v)/2$ lies in $Q$ for all $u,v\in Q$.
This is a condition on the set, not strong convexity of the objective.

\begin{theorem}[A nonzero optimal gradient gives a linear rate]
\label{thm:strong_set_linear}
Suppose additionally that $Q$ is strongly convex and
$\nabla f(x^*)\ne0$. Then $x^*$ is unique and
\[
 c:=\min_{x\in Q}\|\nabla f(x)\|>0.
\]
Define
\[
 \underline\alpha_{\rm sc}=
 \min\left\{1,\left(\frac{L_0}{2\lambda c}+2L_1D^2\right)^{-1}\right\},
\]
with value $1$ when the denominator is zero. Then
\begin{align*}
 f(x_k)-f(x^*)&\le
 (1-\gamma\underline\alpha_{\rm sc}/2)^k
       \bigl(f(x_0)-f(x^*)\bigr),\\
 f(x_k)-\ell_k&\le
 (1-\gamma\underline\alpha_{\rm sc}/2)^k G(x_0)/\gamma.
\end{align*}
The algorithm uses neither $c$ nor $\lambda$; they enter only this rate.
\end{theorem}
\begin{proof}
If $y$ were another minimizer, \eqref{eq:quasar} would make
\[
 \frac{f(x^*+t(y-x^*))-f(x^*)}{t^\gamma}
\]
nondecreasing for $0<t\le1$. It is nonnegative and vanishes at $1$,
so the function is constant on the segment and
$\langle\nabla f(x^*),y-x^*\rangle=0$.
The midpoint ball and constrained first-order optimality instead give
\[
 \langle\nabla f(x^*),y-x^*\rangle
 \ge2\lambda\|\nabla f(x^*)\|\|y-x^*\|^2>0
\]
when $y\ne x^*$. Thus the minimizer is unique. A zero gradient anywhere
in $Q$ would give another minimizer by \eqref{eq:quasar}; continuity
and compactness now give $c>0$.

The same midpoint construction at a full linear minimizer yields
\[
 G(x_k)\ge2\lambda\|\nabla f(x_k)\|\|d_k\|^2
           \ge2\lambda c\|d_k\|^2.
\]
When $L_0=L_1=0$, $\alpha_k=1$. Otherwise,
$\log(1+t)\ge t/(1+t)$ in \eqref{eq:log_step}, or the stated limit
when $L_1=0$, gives
\[
 \alpha_k\ge\min\left\{1,
 \frac{G(x_k)}{(L_0+L_1G(x_k))\|d_k\|^2
                      +L_1D\|d_k\|G(x_k)}\right\}
 \ge\underline\alpha_{\rm sc}.
\]
Here $\|d_k\|\le D$ and the preceding scaling inequality bound the
denominator divided by $G(x_k)$ by $L_0/(2\lambda c)+2L_1D^2$.
The two recurrences follow as in Corollary~\ref{cor:pure_gap_linear}.
\end{proof}
